% !TEX program = xelatex
\documentclass{英语试卷}
\begin{document}
\examtitle{潍坊市 2025 届高三开学调研监测考试}{英语}
\par 潍坊市 2025 届高三开学调研监测考试\par
\par 英语试题\par
\begin{notice}
\par 1. 答卷前，考生务必将自己的姓名、准考证号填写在答题卡上。\par
\par 2. 回答选择题时，选出每小题答案后，用铅笔把答题卡上对应题目的答案标号涂黑。如需改动，用橡皮擦干净后，再选涂其他答案标号。回答非选择题时，将答案写在答题卡上，写在本试卷上无效。\par
\par 3. 考试结束后，将本试卷和答题卡一并交回。\par
\end{notice}
\parttitle{第一部分 听力（共两节，满分 30 分）}
\par 做题时，先将答案标在试卷上，录音内容结束后，你将有两分钟的时间将试卷上的答案转涂到答题卡上。\par
\sectitle{第一节（共5小题：每小题1.5分，满分7.5分）}
\par 听下面 5 段对话。每段对话后有一个小题，从题中所给的 A、B、C 三个选项中选出最佳选项。听完每段对话后，你都有 10 秒钟的时间来回答有关小题和阅读下一小题。每段对话仅读一遍，\par
\question{1. What will the speakers do with the phone?}
\begin{choices}(3)
  \item Have it updated.
  \item Have it charged.
  \item Have it checked.
\end{choices}
\question{2. What kind of T-shirts does the woman prefer?}
\begin{choices}(3)
  \item Short.
  \item Loose.
  \item Tight.
\end{choices}
\question{3. Who is the woman probably?}
\begin{choices}(3)
  \item A language teacher.
  \item A writer.
  \item A musician.
\end{choices}
\question{4. What does the man have with his coffee?}
\begin{choices}(3)
  \item Low-fat milk.
  \item Goat's milk.
  \item Cream.
\end{choices}
\question{5. What are the speakers probably going to do next?}
\begin{choices}(3)
  \item Put up a tent.
  \item Fish in the lake.
  \item Get food at a store
\end{choices}
\sectitle{第二节（共15小题：每小题1.5分，满分22.5分）}
\par 听下面 5 段对话或独白。每段对话或独白后有几个小题，从题中所给的 A、B、C 三个选项中选出最佳选项。听每段对话或独白前，你将有时间阅读各个小题，每小题 5 秒钟；听完后，各小题将给出 5 秒钟的作答时间。每段对话或独白读两遍。\par
\question{6. Where does the conversation probably take place?}
\begin{choices}(3)
  \item In a library.
  \item At a print shop.
  \item In a classroom.
\end{choices}
\question{7. What did the woman do last night?}
\begin{choices}(3)
  \item She shared a story.
  \item She watched a show.
  \item She worked on a presentation.
\end{choices}
\par 听第 7 段材料，回答第 8 至 10 题，\par
\question{8. What does the man plan to do this afternoon?}
\begin{choices}(3)
  \item Go shopping.
  \item Attend a party.
  \item Play tennis.
\end{choices}
\question{9. What does the woman want to do first?}
\begin{choices}(3)
  \item Take a bath.
  \item Buy new clothes
  \item Go into town with the man.
\end{choices}
\question{10. What is the probable relationship between the speakers?}
\begin{choices}(3)
  \item Brother and sister.
  \item Father and daughter
  \item Grandfather and granddaughter
\end{choices}
\par 听第 8 段材料，回答第 11 至 13 题，\par
\question{11. Which month is it now probably?}
\begin{choices}(3)
  \item September.
  \item November.
  \item December.
\end{choices}
\question{12. How does the man feel about coming back to school?}
\begin{choices}(3)
  \item Bored.
  \item Pleased.
  \item Anxious.
\end{choices}
\question{13. What should the man do according to the woman?}
\begin{choices}(3)
  \item Get outside more often.
  \item Take more vitamin pills.
  \item Get more indoor exercise.
\end{choices}
\par 听第 9 段材料，回答第 14 至 17 题。\par
\question{14. Why does the woman come to the man?}
\begin{choices}(3)
  \item To ask for advice.
  \item To have an interview.
  \item To send her congratulations.
\end{choices}
\question{15. Where are the speakers now?}
\begin{choices}(3)
  \item In China.
  \item In the UK.
  \item In Australia.
\end{choices}
\question{16. What are northern British people like?}
\begin{choices}(3)
  \item They have a faster pace of life.
  \item They're more friendly and outgoing.
  \item Their accents are easier to understand.
\end{choices}
\question{17. Which aspect of London attracts the man?}
\begin{choices}(3)
  \item Its cultural factors.
  \item A better career environment.
  \item The comfortable living conditions.
\end{choices}
\par 听第 10 段材料，回答第 18 至 20 题。\par
\question{18. Who did the speaker go to the AI Canteen with?}
\begin{choices}(3)
  \item A family member.
  \item A local friend.
  \item A business partner.
\end{choices}
\question{19. How do people order food in the AI Canteen?}
\begin{choices}(3)
  \item On a screen.
  \item With robotic arms.
  \item Through a waiter.
\end{choices}
\question{20. How many people can the AI system serve in a day?}
\begin{choices}(3)
  \item 100.
  \item 200.
  \item 300.
\end{choices}
\parttitle{第二部分 阅读理解（共两节，满分 50 分）}
\sectitle{第一节（共15小题：每小题2.5分，满分37.5分）}
\par 阅读下列短文，从每题所给的 A、B、C 和 D 四个选项中，选出最佳答案。\par
\passagelabel{A}
\par 2-hour Scooter Rental in Niagara Falls\par
\par If you want to avoid the queues(队列) and crowded areas of visitors, this activity is for you. You can visit scenery that most never do along the Niagara Gorge on an electric scooter(滑板车), which is easy to ride and makes getting from viewpoint to viewpoint fast and easy. You can cover about 3-4 miles of Niagara Falls State Park, the Niagara Gorge, and Whirlpool State Park all in 2 hours using a scooter.\par
\par What's Included\par
\par Ninebot Max Electric Scooter\par
\par Lock to look up the electric scooter\par
\par Meeting and Pickup\par
\par\vspace{0.4em}
\begin{center}\begin{tabular}{{|l|p{\dimexpr\linewidth-3em-4\tabcolsep}|}}
\hline
Meeting point & End point \\
2 Old Falls St.Niagara Falls, NY 14303, USA & This activity ends back at the meeting point. \\
\hline
\end{tabular}\end{center}
\vspace{0.4em}
\par 5/13/2023-7/20/2025 Monday- Sunday: 10:00 AM-06:00 PM\par
\par Additional Information\par
\par Confirmation will be received at time of booking - Not wheelchair accessible - Near public transportation - Not recommended for pregnant(怀孕的) travelers - No heart problems or other serious medical conditions\par
\par Travelers should at least have an average physical fitness level Click here to apply or call: +1 (702) 648-5873 Cancellation Policy\par
\par You can cancel at least 24 hours in advance of the experience for a full refund.\par
\question{21. What is one advantage of renting an electric scooter?}
\begin{choices}(2)
  \item It saves money
  \item It covers all the routes.
  \item It makes your tour easy.
  \item It helps to reduce accidents.
\end{choices}
\question{22. What does the activity offer travellers?}
\begin{choices}(2)
  \item Locks against theft.
  \item Facilities to keep healthy.
  \item Assistance for the disabled.
  \item Guidebooks for Niagara Falls.
\end{choices}
\question{23. When can you call off the activity to get all your money back?}
\begin{choices}(2)
  \item After online confirmation.
  \item One day ahead of the tour.
  \item The moment the tour starts.
  \item Twenty four hours after booking
\end{choices}
\passagelabel{B}
\par Class is in session, and Olivia Munn's 22-month-old son Malcolm is a star student! The mother shared a new video of herself teaching her son Mandarin(普通话) before bed, and he nailed it!\par
\par "Chinese lessons," she wrote over a video of her and Malcolm reading a book before bed. Malcolm stood next to his mom to read the book with her. As Munn said the Mandarin word, the baby repeated it back in the cutest little voice, followed by a smile and a hug for his mama. You can even hear someone say "Aww!" in the background, which was likely Malcolm's dad John Mulaney. Malcolm seemed so proud of himself, and he was so cute!\par
\par In her caption(字幕), Munn explained that the book was a gift from comedian Ali Wong. "Thank you@aliwong for our First 100 Mandarin Words book! Malcolm loves it!" she said.\par
\par This video went viral. Actress Melanie Lynskey commented, "The love!!" Whitney Cummings commented, "My heart." "Growing up multilingual(多语种的) is a magical gift," one person said, "I’ll always be grateful for my mother insisting on speaking Arabic to us \& giving us lessons." Another person commented, "I love this family so much and it’s definitely the cutest family I’ve ever seen."\par
\par The mother often introduces Malcolm to parts of his Asian root. They make noodles for the Lunar New Year, and she's taken him shopping at an Asian market called H Mart. Now he's learning how to speak Mandarin, and it's so adorable.\par
\par In her latest video, Munn is reading My First Mandarin Words with Gordon \& Li Li, written by Michele Wong. The official Gordon and L Li account commented, "This is the sweetest! We're so happy your little guy is having FUN learning Mandarin with Gordon \& Li Li! Big thanks to @ aliwong for sharing the Mandarin love!"\par
\question{24. What does the underlined phrase "nailed it" probably mean?}
\begin{choices}(2)
  \item Got a rich reward.
  \item Hit the nail hard.
  \item Messed things up.
  \item Did a fantastic job.
\end{choices}
\question{25. The second paragraph mainly tells us \blanknum{}.}
\begin{choices}(2)
  \item a bedtime story
  \item an interactive scene
  \item an amusing dad
  \item a reading competition
\end{choices}
\question{26. What can we learn about the video from paragraph 4?}
\begin{choices}(2)
  \item It is well received.
  \item It is full of imagination
  \item It changes the way of parenting.
  \item It inspires people to learn language.
\end{choices}
\question{27. Why is Malcolm exposed to Asian culture?}
\begin{choices}(2)
  \item To learn how to shop
  \item To get to know his origin
  \item To prepare for an Asian trip.
  \item To develop his independence.
\end{choices}
\passagelabel{C}
\par In 2020, there were 727 million people aged 65 or above, which is expected to more than double by 2050. The silver-haired population represent a massive opportunity for nearly every industry where companies can build a new market for their products and services while creating social value.\par
\par The main consumption patterns of elderly population can be classified into two categories: health-related consumption and daily consumption. When it comes to health consumption, seniors are not only concerned about healthcare but also increasingly emphasizes disease prevention. In relation to this, preventive health services, such as physical examinations, will have significant room for growth. Evidently, regular physical examinations require a certain amount of time, economic costs and medical support. Therefore, the high-income population in big cities can sustain them for a longer period of time due to their subjective willingness and objective conditions.\par
\par Another key section of the elderly market is represented by the application of medical devices and technology tailored for seniors. One notable area of development is the adoption of remote monitoring systems. These systems make use of cutting-edge technology to enable healthcare professionals to remotely monitor the vital signs and health status of elderly individuals, providing real-time feedback and intervention when necessary. This not only ensures immediate medical attention but also allows older adults to maintain their independence and age comfortably in their own homes.\par
\par Acknowledging the importance of social engagement and mental stimulation in enhancing the well-being of seniors, the entertainment industry has emerged as a key component as well. It is centered around offering diverse activities and entertainment options, such as organized tours, social clubs, arts and crafts workshops and music and dance performances, which are meant for the unique preferences and needs of older individuals. They offer a platform for them to connect, share experiences and engage, creating a supportive environment where seniors can enrich life experiences and apply their different talents.\par
\par The aging population is one of the most important trends in the world, and it's indeed a powerful force not just socially but economically.\par
\question{28. What is the purpose of the first paragraph?}
\begin{choices}(2)
  \item To raise a topic
  \item To present some figures
  \item To explain a concept.
  \item To provide supportive evidence.
\end{choices}
\question{29. Who is most likely to have a check-up regularly according to paragraph 2?}
\begin{choices}(2)
  \item A jobless worker in Shenzhen.
  \item A fit college student in a big city.
  \item A wealthy old lawyer in Beijing.
  \item A retired poor writer in a rural area.
\end{choices}
\question{30. What benefits do the systems bring to seniors in paragraph 3?}
\begin{choices}(2)
  \item Their health can be tracked distantly.
  \item They can get free medical treatments,
  \item They can receive home visits from doctors.
  \item Their old-age care problems can be settled.
\end{choices}
\question{31. What does the author say about entertainments for seniors?}
\begin{choices}(2)
  \item Limited to music.
  \item Tailored to their needs
  \item Related to their incomes.
  \item Focused on physical challenges.
\end{choices}
\passagelabel{D}
\par New technologies tend to cause a cycle of concern, disorder, and conflict before eventually being accepted. Recent writings about artificial intelligence(AI) and other advances in computer science suggest that we are preparing to welcome the final stage of this latest round of invention.\par
\par The Last Human Job, sociologist Allison Pugh's new book centered on caregiving and human connection in the age of automation, warns readers against unreserved acceptance of these technological advances, citing "connective labor" as valuable human work that will not be easily replaced by algorithms(算法). "Absent from discussions about AI and automation," Pugh argues, "is the impact that these systems might have on the emotional understandings we build of ourselves and others."\par
\par Pugh sets the stage by arguing for the value of human connection and the importance of both seeing and being seen by others. "Practitioners(从业者) said they gained a sense of purpose simply from the opportunity of observing people at their most defenseless," she observes. She then provides a summary of how care work is being increasingly automated, showing how quantification and measurement have taken over many aspects of human-facing jobs, leaving little room for doctors, teachers, and others to build connection.\par
\par "When one goes to a doctor or a teacher, the encounter is full of the potential for shame, a risk that makes it all the more powerful when practitioners show empathic(同理心的) reflection," writes Pugh. In chapter 8, she explains how such interactions can be done right. Research conducted by herself and others identifies three key aspects supporting connective labor: "relational design," or how people are set up to interact with one another; "connective culture," or shared practices and beliefs that influence how people interact; and "resource distribution," which includes time given for interaction, worker-to-client rates, and the extent of technology and data use, among other factors.\par
\par Pugh concludes by arguing that we need a social movement for connection. "We need to fight for what we might call our 'social health'," she maintains.\par
\question{32. What is the main concern of Allison Pugh's new book?}
\begin{choices}(2)
  \item The impact of AI on future employment.
  \item The approach to living and working with AI.
  \item The increasing role of technology in caregiving.
  \item The value of connective labor in the age of automation.
\end{choices}
\question{33. Why does the author mention "practitioners" in paragraph 3?}
\begin{choices}(2)
  \item To explain how to do care work.
  \item To argue against the automation technology.
  \item To prove the significance of human connection.
  \item To discuss the challenges of privacy in caregiving.
\end{choices}
\question{34. What does the underlined word "it" in paragraph 4 refer to?}
\begin{choices}(2)
  \item The encounter.
  \item The shame.
  \item The risk.
  \item The potential.
\end{choices}
\question{35. What is paragraph 4 mainly about?}
\begin{choices}(2)
  \item The overview of chapter 8.
  \item The effective way to interact.
  \item The new trend of caregiving.
  \item The factors of avoiding shame.
\end{choices}
\sectitle{第二节（共5小题；每小题2.5分，满分12.5分）}
\par 阅读下面短文，从短文后的选项中选出能填入空白处的最佳选项。选项中有两项为多余选项。\par
\par In Chinese traditional architecture, a skywell is an open courtyard or well-like space usually located at the heart of a building or group of buildings. \blanknum{36} Their designs often feature the classic Chinese dou gong overhangs to provide shade and collect water into a central basin.\par
\par Chinese skywells are so effective in cooling homes because they can control natural airing. Cool air is drawn into the building from the lower levels, and warm air escapes and rises through the upper levels. \blanknum{37}\par
\par The skywell's layout and construction materials also help in thermal(热量的) regulation. During the day, the courtyard absorbs heat from the sun, acting as a thermal mass. \blanknum{38} This day-night cycle helps maintain a balanced and consistent temperature indoors.\par
\par Most traditional skywells include water features, such as fountains or ponds. A recent study has shown that the most significant part of a skywell's cooling effect comes from these water bodies. \blanknum{39} As water evaporates(蒸发) from these features, it absorbs heat from the surrounding air. This reduces the temperature within the courtyard and surrounding rooms.\par
\par \blanknum{40} One fitting example is the National Heavy Vehicle Engineering Technology Research Centre in Jinan, an 18-story building that has a giant open space at its heart from the fifth floor to the top floor. The toilets, elevators, and meeting rooms surround this space. This enhances natural airing and natural energy consumption, according to the Shanghai-based CCDI Group.\par
\begin{choices}(1)
  \item This promotes fresh air flow, especially during hot summer days.
  \item It is often a square surrounded by the building's rooms and wings.
  \item This cooling was disadvantageous in regions with hot and dry climates
  \item It's probably why temperatures around them are 2.6°C to 4.2°C lower.
  \item These wonders continue to inspire architects to seek sustainable solutions.
  \item The magic of skywells stands as witnesses to the wisdom of ancient builders.
  \item The stored heat is gradually released as the outdoor temperature drops at night
\end{choices}
\parttitle{第三部分 语言知识运用（共两节，满分30分）}
\sectitle{第一节（共15小题；每小题1分，满分15分）}
\par 阅读下面短文，从短文后各题所给的 A、B、C 和 D 四个选项中，选出可以填入空白处的最佳选项。\par
\par As a teenager growing up in Great Britain, Lola Anderson was inspired by the rowing events at the 2012 London Olympics.\par
\par Moved by the athletes' strength and determination, she decided to \blanknum{41} the sport herself. In her diary, she expressed her dream of winning an Olympic gold medal in rowing. Embarrassed by her \blanknum{42} dream, Lola tore out the diary page. "I threw that away because I didn't believe," Anderson \blanknum{43} "I was 14 then, so why would I believe? Young girls struggle to see themselves as strong, athletic individuals, but that's \blanknum{44} now." Despite her initial \blanknum{45}, Anderson pursued rowing with her father's support.\par
\par In 2019, as Don Anderson \blanknum{46} cancer, he presented Lola with a \blanknum{47}. He held open his hand to \blanknum{48} the page she had torn from her diary years earlier. Don had found it in the trash(垃圾筒) and kept it, \blanknum{49} she would need it one day. Don \blanknum{50} months later, but his faith in his daughter's dream remained.\par
\par His predictive gesture \blanknum{51} on Wednesday when Anderson competed in her first Olympic Games as part of the women's quadruple sculls(四人双桨) rowing team. Her team \blanknum{52} the gold medal by a mere 0.15 seconds. After the race, Anderson reflected on her father's firm support.\par
\par "It's a piece of paper, but it's the most valuable thing I have," she said. "Maybe jointly with the \blanknum{53} now."\par
\par Lola Anderson's 54 from a self-doubting teenager to an Olympic champion serves as 55 of the\par
\par power of dreams and the lasting impact of a father's love.\par
\cz{41}{take up}{look into}{live upon}{fight for}
\cz{42}{greedy}{achievable}{shallow}{wild}
\cz{43}{imagined}{regretted}{added}{recalled}
\cz{44}{changing}{strengthening}{worsening}{speeding}
\cz{45}{resolution}{doubt}{confidence}{worry}
\cz{46}{studied}{battled}{defeated}{prevented}
\cz{47}{wish}{blow}{promise}{surprise}
\cz{48}{throw}{reveal}{fold}{release}
\cz{49}{advocating}{proving}{feeling}{wondering}
\cz{50}{passed away}{died off}{set off}{went away}
\cz{51}{happened}{mattered}{arrived}{worked}
\cz{52}{bagged}{forgotten}{lost}{recovered}
\cz{53}{support}{team}{medal}{rowing}
\cz{54}{journey}{range}{departure}{achievement}
\cz{55}{advice}{belief}{memory}{proof}
\sectitle{第二节（共10小题；每小题1.5分，满分15分）}
\par 阅读下面短文，在空白处填入 1 个适当的单词或括号内单词的正确形式，\par
\par A tea culture salon may be the best way to experience China's tea ceremonies and culture in overseas locations.\par
\par "Tea for Harmony" cultural salon in Bath, England on Friday was aimed \blanknum{56} both celebrating International Tea Day, which fell on May 21 and promising participants an opportunity \blanknum{57} (experience) teas and culture from Central China's Hubei province. The \blanknum{58} (represent) group from the Hubei Tea Association brought three types of tea to be enjoyed at the salon: Enshi Yulu, Yihong Black Tea, and Qingzhuan Tea\par
\par The history of the teas on offer \blanknum{59} (date) back more than 1,000 years, and the crafting techniques of Enshi Yulu and Qingzhuan Tea have become national intangible cultural heritages. Each tea was presented with its \blanknum{60} (appeal) tea ceremony, where a tea master introduced its area of origin and its distinct characteristics, \blanknum{61} demonstrated the particular way the tea should be prepared and appreciated. "Unique", "delicate" and "wonderful" \blanknum{62} (be) words British participants used to describe their feelings during the tea tastings.\par
\par "Tea is such an important part of British culture that we have an entire meal \blanknum{63} (create) that we have entirely dedicated to our afternoon tea. It's very important for us British to be in Bath in this afternoon, \blanknum{64} is the best and most elegant afternoon tea destination," said Dianne Francombe, CEO of the Bristol and West of England China Bureau.\par
\par "Fine tea is more expressive than fine wine, and that experience, the flavors and delicacy of tea, is something \blanknum{65} (true) superb. I look forward to continued friendship with friends from China whenever you visit this region."\par
\parttitle{第四部分 写作（共两节，满分 40 分）}
\sectitle{第一节（满分15分）}
\par 假如你是李华，你校英文报举办了主题为“终身学习(Lifelong Learning)”的读书分享会活动。请你写一篇投稿记录这次活动，内容包括：\par
\par （1）活动的内容；(2) 你的心得。\par
\par\vspace{0.4em}
\begin{center}\begin{tabular}{{|l|}}
\hline
English Song Appreciation Course \\
 \\
 \\
\hline
\end{tabular}\end{center}
\vspace{0.4em}
\par 注意:\par
\par (1) 写作词数应为 80 左右;\par
\par (2) 请按如下格式在答题纸的相应位置作答。\par
\sectitle{第二节（满分25分）}
\par 阅读下面材料，根据其内容和所给段落开头语续写两段，使之构成一篇完整的短文。\par
\par Wil is was a highway patrol trooper(巡逻警察), who was in charge of the traffic safety of the Tennessee Highway. One afternoon this July, it rained heavily. He was driving along the busiest section of the highway during the downpour when he saw the traffic slowing down in front of him. Some cars stopped and bypassed one lane(车道).\par
\par Willis initially thought there was a car crash, so he slowed down his car to approach. But he soon found out the source of the delay was an animal wandering on the wet road. And he spotted that was a dog when he got close enough. There was a large dog that looked like a brown German shepherd(牧羊犬) in the middle lane\par
\par The weather was at its worst, and the gray sky with the heavy rain gave the highway minimum visibility. The delay might cause a traffic jam and of course it was very dangerous to let the dog wandering in the middle of the highway.\par
\par He immediately turned on his blue lights because people usually slow down when they see the blue lights and called into the radio so they knew where he was. He pulled over, got out, and whistled to draw the dog's attention.\par
\par It seemed as if the dog knew Willis was its lifesaver. It walked to Willis and his car. Once the dog was in the car, Willis pulled into a parking lot nearby, where he called Networks Animal Service. While waiting, he noticed the dog wearing a collar with its name "Luna". Willis recorded a short video of the dog.\par
\par Then an animal control officer came and took Luna to the shelter. Willis came back to work. Watching cars travelling on the highway through the window glass, he thought he might do something to let the dog's owner know their pet was found safe.\par
\par 注意：\par
\par (1) 续写词数应为 150 左右;\par
\par (2) 请按如下格式在答题卡的相应位置作答。\par
\writingpara{Paragraph 1: After work, Willis posted the video online with a brief description attached.}
\writinglines{8}
\writingpara{Paragraph 2: Days later, news came that Luna's owner had been found and would visit him.}
\writinglines{8}
\end{document}
